\documentclass[11pt,a4paper]{article}

\usepackage[utf8]{inputenc}
\usepackage[french]{babel}
\usepackage[T1]{fontenc}
\usepackage{graphicx}
\usepackage{xcolor}
\usepackage{hyperref}
\usepackage{geometry}
\usepackage{titlesec}
\usepackage{enumitem}
\usepackage{tcolorbox}
\usepackage{fancyhdr}
\usepackage{bookmark}

\geometry{margin=2cm}

% --- COULEURS WATERPAID ---
\definecolor{wpblue}{HTML}{0061FF}
\definecolor{wpcyan}{HTML}{00D1FF}
\definecolor{wpdark}{HTML}{0F172A}
\definecolor{wpbg}{HTML}{F8FAFC}
\definecolor{success}{HTML}{22C55E}
\definecolor{warning}{HTML}{F59E0B}

% --- CONFIGURATION DES TITRES ---
\titleformat{\section}{\color{wpblue}\normalfont\Large\bfseries}{\thesection.}{0.5em}{}[\titlerule]
\titleformat{\subsection}{\color{wpblue}\normalfont\large\bfseries}{\thesubsection.}{0.5em}{}

% --- CONFIGURATION DU HEADER/FOOTER ---
\pagestyle{fancy}
\fancyhf{}
\lhead{\color{wpblue}\textbf{WaterPaid} - Manuel Complet}
\rhead{\thepage}
\lfoot{\small © 2026 WaterPaid - L'eau Intelligente}

% --- PAGE DE GARDE ---
\title{
    \vspace{2cm}
    \textbf{\Huge WaterPaid} \\
    \vspace{0.5cm}
    {\large \textit{Guide d'Installation et d'Utilisation}} \\
    \vspace{1.5cm}
    \includegraphics[width=0.3\textwidth]{logo.png} \\ 
    \vspace{1cm}
    \textbf{\LARGE MANUEL UTILISATEUR COMPLET} \\
    \vspace{0.5cm}
    \textit{Pour les Gestionnaires et les Familles} \\
    \vspace{1cm}
    \small{Version 4.0 — Édition Exhaustive}
}
\author{Documentation WaterPaid}
\date{\today}

\begin{document}

\maketitle
\thispagestyle{empty}
\newpage

\tableofcontents
\newpage

\section{Introduction}
Ce guide est conçu pour vous accompagner pas-à-pas dans l'utilisation de l'écosystème **WaterPaid**. 
Nous avons séparé les instructions en deux grandes parties : 
\begin{enumerate}
    \item **Pour le Gestionnaire (Admin)** : De la réception du matériel à la gestion quotidienne.
    \item **Pour l'Utilisateur (Client)** : De l'installation de l'application à l'achat d'eau.
\end{enumerate}

---

\section{PARTIE 1 : LE GESTIONNAIRE (ADMINISTRATEUR)}
\textit{Vous êtes le responsable technique ou le distributeur qui installe les compteurs.}

\subsection{Phase 1 : Mise en service du matériel (L'Atelier)}
Avant même d'aller chez le client, vous devez préparer les compteurs.

\subsubsection{1. Allumage du compteur}
Chaque compteur possède un interrupteur de sécurité pour le transport.
\begin{itemize}
    \item Localisez le petit interrupteur sur le côté ou en-dessous du boîtier.
    \item Basculez-le sur la position **"AL"** (Allumé).
    \item Le compteur va émettre un signal ou un voyant va clignoter pour indiquer qu'il est vivant.
\end{itemize}

\begin{figure}[h]
    \centering
    \includegraphics[width=0.4\textwidth]{images/hardware_switch.png}
    \caption{Mise sous tension : Interrupteur sur "AL"}
\end{figure}

\subsubsection{2. Création de votre compte Admin}
Maintenant que le compteur est allumé, il faut le dire à l'application.
\begin{enumerate}
    \item Téléchargez l'application **WaterPaid**.
    \item Cliquez sur **"S'inscrire"**.
    \item **IMPORTANT** : Sélectionnez le rôle **"ADMINISTRATEUR"** (et non Utilisateur).
    \item Entrez votre nom, votre numéro pro et un mot de passe sécurisé.
\end{enumerate}

\begin{figure}[h]
    \centering
    \includegraphics[width=0.4\textwidth]{images/inscription.png}
    \caption{Écran d'inscription (Bien choisir Admin)}
\end{figure}

\subsubsection{3. Enregistrement du compteur dans le système}
C'est l'étape clé où le compteur devient "intelligent".
\begin{enumerate}
    \item Connectez-vous à votre compte Admin.
    \item Allez dans l'onglet **Compteurs** (icône Jauge).
    \item Cliquez sur le bouton **"Ajouter"**.
    \item Remplissez les champs :
        \begin{itemize}
            \item **Numéro de Série (S/N)** : Écrit sur l'étiquette du compteur.
            \item **DevEUI** : L'identifiant radio unique (souvent un code long avec des chiffres et lettres).
        \end{itemize}
    \item Validez. L'application va vous générer un **TOKEN** (un code secret unique).
    \item **NOTEZ CE TOKEN** : Vous devrez le donner au client final.
\end{enumerate}

\begin{figure}[h]
    \centering
    \includegraphics[width=0.4\textwidth]{images/admin_create_meter.png}
    \caption{Formulaire d'ajout de compteur}
\end{figure}

\subsection{Phase 2 : L'Installation chez le Client (Le Terrain)}
Tout est prêt numériquement. Passons à l'installation physique.

\begin{enumerate}
    \item **Sécurité** : Éteignez le compteur (Interrupteur sur OFF) pour le transport jusqu'au domicile.
    \item **Plomberie** : Coupez l'arrivée d'eau principale.
    \item Installez le compteur **juste après** le robinet d'arrêt général. Assurez-vous que les flèches sur le tuyau suivent le sens de l'eau.
    \item **Mise en marche finale** : Une fois posé, remettez l'interrupteur sur **"AL"**.
    \item Ouvrez l'eau doucement. Vérifiez qu'il n'y a pas de fuite.
\end{enumerate}

\subsection{Phase 3 : Vos Outils de Gestion Quotidienne}
Votre application Admin est un véritable tableau de bord professionnel.

\subsubsection{Le Dashboard (Vue d'ensemble)}
Dès l'ouverture, vous voyez la santé de votre réseau :
\begin{itemize}
    \item **Revenus Totaux (Vert)** : L'argent encaissé par toutes les recharges.
    \item **Volume (Bleu)** : Le total des mètres cubes d'eau distribués.
    \item **Compteurs Actifs (Orange)** : Nombre d'appareils qui répondent bien au réseau.
    \item **Clients (Gris)** : Nombre de familles desservies.
\end{itemize}

\begin{figure}[h]
    \centering
    \includegraphics[width=0.6\textwidth]{images/dashboard_admin.png}
    \caption{Tableau de bord Administrateur}
\end{figure}

\subsubsection{Gérer les Utilisateurs et les Problèmes}
Dans l'onglet **Utilisateurs**, vous pouvez aider un client en difficulté :
\begin{itemize}
    \item **Recherche** : Trouvez un client par son nom ou son numéro.
    \item **Détails** : Voyez s'il a du crédit ou si sa vanne est fermée.
    \item **Recharge Manuelle** : Si un client vous paye en espèces (CASH), c'est ici que vous lui envoyez l'eau manuellement.
    \item **Rapports** : Dans les paramètres, téléchargez un bilan PDF de vos activités.
\end{itemize}

---

\section{PARTIE 2 : L'UTILISATEUR (LE FOYER)}
\textit{Vous êtes le client qui consomme l'eau à la maison.}

\subsection{Phase 1 : Bienvenue chez WaterPaid}
Votre compteur a été installé par le pro. Voici comment prendre le contrôle.

\subsubsection{1. Créer son compte}
\begin{enumerate}
    \item Ouvrez l'application **WaterPaid**.
    \item Cliquez sur **"S'inscrire"**.
    \item Choisissez le rôle **"UTILISATEUR"** (c'est le choix par défaut).
    \item Mettez votre nom et votre numéro de téléphone. C'est tout !
\end{enumerate}

\subsubsection{2. Lier votre compteur (La clé magique)}
Pour voir votre consommation, il faut dire à l'application quel est *votre* compteur.
\begin{enumerate}
    \item Sur l'accueil, cliquez sur le gros bouton **"+" Bleu**.
    \item **Option A (Facile)** : Scannez le QR Code que l'installateur a collé sur le compteur.
    \item **Option B (Manuelle)** : Tapez le **Token** (le code secret) que l'installateur vous a donné.
    \item Cliquez sur **"Lier"**. Une carte blanche apparaît : c'est votre compteur !
\end{enumerate}

\begin{figure}[h]
    \centering
    \includegraphics[width=0.4\textwidth]{images/scan_qr.png}
    \caption{Liaison du compteur par QR Code ou Token}
\end{figure}

\subsection{Phase 2 : Utilisation Quotidienne}

\subsubsection{Le Tableau de Bord (Dashboard)}
Tout ce qu'il faut savoir en un coup d'œil :
\begin{itemize}
    \item **Solde Global (En haut)** : L'argent total disponible dans votre portefeuille virtuel.
    \item **Vos Compteurs** : Si vous en avez plusieurs (Maison, Boutique), ils sont listés ici.
    \item **Dernières activités** : Pour savoir quand vous avez rechargé la dernière fois.
\end{itemize}

\begin{figure}[h]
    \centering
    \includegraphics[width=0.4\textwidth]{images/dashboard_user.png}
    \caption{Votre espace personnel}
\end{figure}

\subsubsection{Tout savoir sur votre compteur (Les Détails)}
Cliquez sur votre compteur pour voir ses secrets.
\begin{itemize}
    \item **Jauge d'Eau** : Elle se vide au fur et à mesure que vous consommez.
    \item **Icône Signal (Radio)** : Indique si le compteur capte bien le réseau.
    \item **Icône Pile (Batterie)** : Indique si la batterie du compteur est bonne.
    \item **La Vanne (Robinet)** : 
        \begin{itemize}
            \item \textbf{Cercle Vert} = Ouverte, l'eau coule.
            \item \textbf{Cercle Rouge} = Fermée, pas de crédit ou problème.
        \end{itemize}
\end{itemize}

\begin{figure}[h]
    \centering
    \includegraphics[width=0.4\textwidth]{images/details_compteur.png}
    \caption{Détails techniques simplifiés}
\end{figure}

\subsubsection{Comment acheter de l'eau ? (Recharger)}
Plus d'eau ? Pas de panique, ça prend 30 secondes.
\begin{enumerate}
    \item Cliquez sur le bouton **"Recharger"**.
    \item Tapez le montant en FCFA (ex: 1000 F).
    \item Choisissez **Mobile Money** (Orange/MTN) ou **Cash** (si déjà payé).
    \item Validez sur votre téléphone.
    \item **Magie !** En quelques instants, le compteur reçoit l'ordre et l'eau revient.
\end{enumerate}

\begin{figure}[h]
    \centering
    \includegraphics[width=0.4\textwidth]{images/recharge_payunit.png}
    \caption{Recharge facile par Mobile Money}
\end{figure}

\subsection{Phase 3 : Vos Paramètres}
Dans l'onglet **Profil** (icône Bonhomme), vous pouvez :
\begin{itemize}
    \item Changer votre mot de passe.
    \item Passer l'application en **Mode Sombre** (Dark Mode) pour ne pas vous éblouir la nuit.
    \item Contacter le **Support** si vous avez une fuite ou un souci.
\end{itemize}

\begin{figure}[h]
    \centering
    \includegraphics[width=0.4\textwidth]{images/profil_settings.png}
    \caption{Réglages et préférences}
\end{figure}

---

\section{Dépannage Rapide (FAQ)}

\textbf{Q: J'ai payé mais l'eau ne revient pas tout de suite.} \\
R: Le réseau radio peut prendre jusqu'à 2 minutes pour recevoir l'ordre. Attendez un peu. Si ça dépasse 5 minutes, contactez le support.

\textbf{Q: Je n'arrive pas à scanner le QR Code.} \\
R: Vérifiez que vous avez autorisé l'application à utiliser votre caméra. Sinon, entrez le code manuellement.

\textbf{Q: La batterie du compteur est faible (Rouge).} \\
R: Prévenez votre gestionnaire. Il viendra changer la pile. N'essayez pas de l'ouvrir vous-même !

\vspace{2cm}
\begin{center}
    \textit{Merci d'utiliser WaterPaid. L'eau potable, gérée intelligemment.}
\end{center}

\end{document}
